\documentclass[11pt,a4paper]{article}
\usepackage[utf8]{inputenc}
\usepackage{cmap}
\usepackage[T1]{fontenc}
\usepackage{lmodern}
\usepackage[spanish,es-noquoting]{babel}
\usepackage[margin=2.0cm,top=1.6cm,bottom=1.7cm]{geometry}
\usepackage{booktabs}
\usepackage{parskip}
\usepackage{amsmath}
\usepackage{caption}
\renewcommand{\arraystretch}{1.15}
\captionsetup{font=small,labelfont=bf}
\setlength{\parskip}{0.38em}

\begin{document}

\begin{center}
{\large\bfseries Informe de cierre --- frente de archivos cifrados}\\[6pt]
{\small 28 de septiembre de 2026 \quad\textbullet\quad Para: Prof.\ Cristian Cappo \quad\textbullet\quad
De: Romina Alfonzo y Carlos Urdapilleta}\\[4pt]
{\footnotesize\itshape Detección de familias de ransomware en base a archivos encriptados
y notas de rescate}
\end{center}

\vspace{0.3em}
\hrule
\vspace{0.5em}

El frente de archivos cifrados queda cerrado: los dos elementos de acción de la reunión del
12/08 están medidos, se agregaron experimentos que llevan el resultado publicado de 0,912 a
0,9998 de exactitud, la cifra final pasó la misma validación que la anterior, y todo está
redactado en el Capítulo~4. Este informe resume qué quedó, qué se probó y qué sigue abierto. El
detalle completo, con bases y semillas, está en el informe del 27/09 (actualizado el 29/09) y en
\texttt{ESTADO\_TESIS.md}.

\section*{1. La cifra del frente}

\begin{center}
\small
\begin{tabular}{lccl}
\toprule
\textbf{Configuración} & \textbf{Exactitud} & \textbf{Macro-F1} & \textbf{Alcance} \\
\midrule
Bytes 512+512 (Exp.\ 2c) & 0,912 $\pm$ 0,002 & 0,911 $\pm$ 0,001 & base del sistema \\
+ rasgos estructurales (Exp.\ 2e) & 0,9357 $\pm$ 0,0005 & 0,9359 $\pm$ 0,0004 & solo contenido \\
\textbf{+ forma de la extensión (Exp.\ 2g)} & \textbf{0,9998 $\pm$ 0,0001} & \textbf{0,9998 $\pm$ 0,0001} & \textbf{canónica} \\
\bottomrule
\end{tabular}
\end{center}

Todas sobre 15.000 archivos y 30 familias. La cifra del frente es la del sistema que apila las
tres capas: las treinta familias superan 0,99 de F1, y bajo tipo de documento no visto da 0,9983 $\pm$ 0,0011
de macro-F1 (cinco semillas). Se reporta siempre con su descomposición: la parte de contenido (0,936 de macro-F1) no
emplea el nombre; la capa de la extensión la eleva a 0,9998 y es la más expuesta a la limitación de campaña.

\section*{2. Qué se probó}

\textbf{Aporte del nombre (elemento de acción 2).} La extensión literal, sola, alcanza 0,9244 de
macro-F1 ---más que los bytes--- y se declara como cota superior: 25 de 30 familias tienen una
única extensión. La forma del nombre, sola, alcanza apenas 0,5771; con los bytes, 0,9998: la
mejora es de la combinación.

\textbf{Cuántos archivos hacen falta.} Con 200 archivos por familia el frente está en 0,9045 de
macro-F1; los 300 restantes hasta 500 aportan +0,0072 en total.

\textbf{Rasgos estructurales (Exp.\ 2e).} Cuarenta y cuatro rasgos que describen la forma del
archivo y no su contenido. Macro-F1 de 0,9114 a \textbf{0,9359}: $\Delta$ = +0,0246 [+0,0237;
+0,0254], 5/5 semillas, concentrado en las seis familias que el 2c no resolvía. Lo sostiene el
tamaño en combinación con el perfil de entropía; el rasgo individual más informativo es el resto
del tamaño módulo 16, huella del modo de cifrado. Bajo tipo no visto, bytes y bytes + estructura
pierden lo mismo (0,049 y 0,050) y el delta se conserva (+0,0238 [+0,0108; +0,0368], 7/7, sin el caso degenerado de BLACKMATTER, sección 3).

\textbf{El sistema completo (Exps.\ 2f y 2g).} Las tres capas apiladas dan 0,9998 en validación
cruzada, pero con la forma del nombre \emph{completo} la validación por tipos colapsó en
\texttt{jpg}: de 0,8052 a 0,2167. Causa, verificada mirando los archivos: la base del nombre la
puso NapierOne, y en las imágenes lleva un \texttt{-fromweb} (\texttt{0001-jpg-fromweb.jpg.avos2}
frente a \texttt{0001-doc.doc.avos2}); los rasgos aprendieron también eso. Con catorce rasgos que
no miran la base del nombre ---trece sobre la extensión final--- se mantiene el 0,9998 y el colapso desaparece: \texttt{jpg} 1,0000,
promedio 0,9963, $\Delta$ = +0,1149 [+0,0743; +0,1554] sobre bytes + estructura, 7/7; con cinco semillas, 0,9983 $\pm$ 0,0011.
Al pasar a un tipo no visto pierde 0,0015, frente a unas cinco centésimas de las capas de contenido.

\textbf{Qué aporta cada capa (Exp.\ 2h).} La forma de la extensión sola da 0,878 de macro-F1: no
reemplaza al contenido. Con ella, cualquiera de las dos capas de contenido basta (estructura +
extensión, 58 rasgos: 0,9997; bytes + extensión: 0,9999): la extensión separa a las familias que
el contenido confunde, y el contenido a las que comparten la forma de su extensión.

\section*{3. Qué se encontró en los datos}

\textbf{CERBER cifra parcialmente.} Los 988 archivos de la familia conservan la cabecera del
documento original (verificado por firma en los 988; en una muestra de seis, entropía de cabecera
0,88 a 6,46 frente a 7,58 a 7,64 en cuerpo y cola). Es un comportamiento del programa y se
declara como resultado sobre la familia.

\textbf{CRYPTOLOCKER cifra de forma determinista.} Los 143 \texttt{.doc} de la familia empiezan
con los mismos 16 bytes cifrados, y los 143 \texttt{.xls} también: la huella de una clave fija sin
vector de inicialización por archivo. Es una regularidad de contenido por tipo de documento, que el
criterio del Exp.~2b no podía ver, y depende de la clave, que es de la campaña.

\textbf{Un sesgo en la carga, corregido.} El filtro que excluía la documentación de NapierOne
descartaba 310 muestras cifradas de BADRABBIT y NOTPETYA. Corregido, la exactitud es 0,9123,
frente a 0,9128 sin corregir y 0,912 publicado: indistinguibles. Se apartaron además 40 archivos
sin cifrar (0,13\,\%), y quedan declarados dos PDF de JIGSAW con cabecera en claro.

\textbf{Censo y duplicados (Exp.\ 2h).} BLACKMATTER son solo imágenes (988 de 1.000): sin ellas
queda con 7 muestras de entrenamiento, y la primera validación por tipos, corrida sin la ponderación
de clases de la publicada, no la reconocía; con el modelo correcto solo cambia ese pliegue. Tres
documentos repetidos en la base de NapierOne dan 16 pares de archivos cifrados idénticos en seis
familias, que cifran de forma determinista; ninguno entre familias distintas.

\section*{4. Qué no salió como se predijo}

Cada experimento llevó sus predicciones escritas y commiteadas antes de correr. Fallaron nueve:
la forma del nombre sola daría 0,97 (dio 0,58); los rasgos estructurales solos darían 0,40 a
0,60 (dieron 0,87); las que más mejorarían serían SUNCRYPT y NOTPETYA (fueron WASTEDLOCKER y
JIGSAW); \texttt{pdf} sería de los tipos que peor generalizan (fue el segundo mejor); solo bytes
promediaría 0,86 a 0,90 bajo tipo no visto (dio 0,85 sin la ponderación de clases de la medición publicada; con ella, 0,876); y la forma del
nombre mejoraría en los siete pliegues de tipo no visto (colapsó en \texttt{jpg}). Las cuatro
predicciones del Exp.~2g, que corrigió ese colapso, se cumplieron. En el Exp.~2h fallaron tres de
dieciséis: duplicados en seis familias y no en dos; la estructura por debajo de los bytes en
\texttt{jpg}, por BLACKMATTER con 7 ejemplos; y el sistema completo en 0,980 y no 0,99 en el pliegue
\texttt{pdf} de una semilla (umbral mal fijado).

\section*{5. Qué queda abierto}

\begin{itemize}
\item \textbf{Una campaña por familia.} Ninguna medición sobre NapierOne separa «identifica la
familia» de «identifica esa campaña»; está definido en metodología y declarado en cada resultado.
La capa de la extensión es la más expuesta: aprende el esquema de renombrado de cada campaña.
Requiere otro conjunto de datos.
\item \textbf{Las familias difíciles, por contenido.} Bajo solo contenido, NOTPETYA, JIGSAW y
CRYPTOLOCKER quedan por debajo de 0,75 de F1 (promedio de cinco semillas; DARKSIDE en 0,7515). El
sistema completo las lleva por encima de 0,99, con la limitación de campaña pegada.
\item \textbf{Combinación con el frente de notas (elemento de acción 3).} Con el sistema completo
no hay margen: ninguna combinación puede sumar más de 0,0002 de exactitud. Sin el nombre del
archivo, delegar en las notas ganaría $\approx$ +2,5 puntos (techo del oráculo: +5,2); no es evaluable
sin muestras pareadas, que no existen.
\end{itemize}

\section*{6. Estado del documento}

Capítulo~4 con los Experimentos 2d a 2h, el censo de integridad, el sesgo corregido y la curva de
aprendizaje; síntesis, tabla comparativa y definición de campaña actualizadas. Compila: 125 páginas, con
el frente de notas integrado, sin errores. Pendiente de su conformidad: 0,9998 del sistema completo como
cifra del frente, con 0,936 de contenido y 0,912 de solo bytes al lado.

\end{document}
